\documentclass{report}
\usepackage{amsmath,amssymb}
\setlength{\parindent}{0mm}
\setlength{\parskip}{1em}
\begin{document}
\begin{center}
\rule{6in}{1pt} \
{\large Jie Chen \\
{\bf Scalable Statistical Analysis of Gaussian Models for Petascale Spatiotemporal Data}}

Mathematics and Computer Science Division \\ Argonne National Laboratory \\ 9700 S Cass Ave Bldg 240 \\ Argonne \\ IL 60439
\\
{\tt jiechen@mcs.anl.gov}\\
Mihai Anitescu\\
Michael Stein\end{center}

Gaussian processes are widely used throughout the statistical and machine
learning communities for modeling natural processes with broad
applications to fields such as nuclear engineering and climate science.
Such models typically have some vector of unknown parameters that must be
estimated to carry out further data analysis tasks such as regression,
classification and interpolation. A standard statistical approach to
estimating these parameters is to choose the parameter values that
maximum the likelihood function (the joint density of the observations).
Simple in principle, the technique is difficult to apply in practice in
the case of very large data sets due to the need to work with the
covariance matrices of the observations. In some cases, the covariance
matrices (or their inverse) may have some exploitable properties
(sparseness, Toeplitz) to reduce computations and/or storage, but in many
applications, the covariance matrices are dense and unstructured. The
likelihood function for Gaussian processes involves a quadratic form in
the inverse covariance matrix and the log determinant of the covariance
matrix. Existing algorithms for maximizing the likelihood heavily rely on
the Cholesky factorization, the computation of which is prohibitively
costly for many problems of practical interest. We propose a sample
average formulation of the optimization framework, which narrows down the
several linear algebra challenges to solving a linear system with the
covariance matrix for multiple right-hand sides. We further investigate
two of the most important ingredients in the conjugate gradient solver:
the conditioning of the covariance matrix and the multiplication of a
vector to a covariance matrix. These two aspects require different
designs of the preconditioner and the matrix-vector multiplication in
different grid configurations. We demonstrate the successful scalable
resolution of the maximum likelihood problem for data sizes as large as a
million points on a grid on a single desktop machine, whereas the
Cholesky factorization approach would have needed a moderate-size
supercomputer. In addition, we have proved that in some circumstances
optimal preconditioning is achievable by means of a filtering approach.
Parallel programs are under development to solve the problem for much
larger data sets and in higher dimensions.


\end{document}
